\begin{summarybox}
	\textbf{\textcolor{CSummary}{一句话核心}}
	\textbf{利用角平分线的对称性，通过垂线、截等长、作平行三种辅助线构造全等或等腰三角形；等腰三角形中顶角平分线、中线、高互相重合。}

	\medskip
	\textbf{\textcolor{CSummary}{使用条件}}
	\begin{itemize}[leftmargin=*, itemsep=0.5em, topsep=0.4em]
		\item \textbf{条件 1（角平分线存在）：} 图形中出现角平分线，或隐含角平分条件。
		\item \textbf{条件 2（点在角分线上）：} 所作垂线、截等长或平行线的起点必须在角平分线上。
		\item \textbf{条件 3（等腰前提）：} 三线合一模型需要三角形为等腰三角形。
	\end{itemize}

	\medskip
	\textbf{\textcolor{CSummary}{关键提醒}}
	\begin{itemize}[leftmargin=*, itemsep=0.5em, topsep=0.4em]
		\item \textbf{易错点（垂线段误用）：} 角平分线性质仅保证垂线段相等，斜线段不一定相等。
		\item \textbf{检查项（模型选用）：} 根据题目条件选择匹配的辅助线模型，不可混用条件。
	\end{itemize}

\end{summarybox}
